\chapter{Conclusiones}
\label{chap:conclusiones}

\noindent
Este trabajo evaluó si el aumento de la fiscalización vial asociado a la
renovación de la Policía de Tránsito de la Ciudad de México (abril de 2022)
redujo las víctimas de hechos de tránsito, comparando la evolución de las
colonias expuestas a la infraestructura de fiscalización con la de las no
expuestas.

No se encuentra evidencia robusta de un efecto. Las víctimas fallecidas —el
\emph{outcome} menos sensible a los efectos de reporte— no cambian en ninguna
especificación. El aumento aparente en las víctimas totales bajo mínimos
cuadrados no sobrevive a un modelo de conteo y coincide en momento y duración
con un repunte transitorio de las infracciones, lo que sugiere un efecto de
registro más que un cambio real en la siniestralidad.

\section{Cómo leer el nulo}

Un resultado nulo admite al menos dos lecturas, y conviene ser explícito sobre
cuál corresponde a este trabajo. La primera es que la fiscalización vial, bien
implementada, no tiene capacidad de reducir víctimas de tránsito —un hallazgo
fuerte, y de alguna forma sorprendente a la luz de \citet{deangelohansen2014},
quien sí encuentra un efecto grande del enforcement policial sobre la
mortalidad vial en su contexto—. La segunda, más modesta y más consistente
con la evidencia reunida aquí, es que \emph{esta} intervención en particular
—la renovación de la Policía de Tránsito— no generó una fiscalización
diferencial lo bastante fuerte ni lo bastante sostenida como para mover el
outcome, independientemente de si el mecanismo funcionaría con una
intervención de mayor intensidad. El mecanismo (Sección~\ref{sec:mecanismo})
muestra un incremento de apenas 10 a 35\,\% en infracciones cerca de la
infraestructura, nunca estadísticamente significativo de forma individual, y
concentrado en una ventana de aproximadamente cuatro meses antes de disiparse.
Un primer paso tan débil deja poco margen para observar un efecto en el
segundo. El nulo de esta tesis debe leerse, entonces, como \emph{ausencia de
evidencia de un efecto bajo un tratamiento débil}, y no como una estimación
precisa de un efecto verdaderamente cero. La réplica en septiembre de 2023
(Sección~\ref{sec:replica_sep2023}) refuerza la robustez del patrón —el mismo
nulo aparece en un segundo evento, con otra reforma y otra infraestructura de
referencia— pero no resuelve esta ambigüedad de fondo, porque tampoco en ese
caso hay evidencia de una fiscalización diferencial fuerte.

\section{Relación con la versión previa de esta investigación}

La versión anterior de este proyecto estimaba, con un diseño de variables
instrumentales a nivel código postal, que un mayor volumen de multas se
asociaba con \emph{más} accidentes de tránsito. El instrumento —haber
recibido infracciones por encima de la mediana en el periodo previo— no
satisfacía de forma creíble la restricción de exclusión: es difícil sostener
que el historial de multas de una zona afecta la accidentalidad únicamente a
través de las multas contemporáneas, y no a través de características de la
zona (tipo de vialidad, volumen de tráfico, comportamiento de los
conductores) correlacionadas con ambas variables. El rediseño que da lugar a
esta tesis corrige tres decisiones de ese diseño original: cambia la unidad de
análisis de código postal a colonia, que es la unidad administrativa sobre la
que efectivamente varía la infraestructura de fiscalización; abandona el
enfoque de variables instrumentales en favor de un estudio de evento con
efectos fijos, sin pretender resolver la endogeneidad de las multas por la vía
de un instrumento débil; y trata las multas explícitamente como
\emph{mecanismo} —evidencia de que hubo más fiscalización— en vez de como la
variable de interés central, desplazando esa posición a las víctimas de
tránsito. El resultado cambia en consecuencia: de un efecto positivo y
significativo, pero apoyado en un instrumento cuestionable, a un nulo robusto
a la especificación, apoyado en una fuente de variación más transparente
—dónde se ubica la infraestructura de fiscalización— aunque también más
débil.

\section{Implicaciones y trabajo futuro}

Con las salvedades del nulo señaladas arriba, este trabajo deja tres
implicaciones. Primero, una reforma institucional de la policía de tránsito
—cambiar quién puede infraccionar y cómo, sin necesariamente ampliar de forma
sostenida la infraestructura física de fiscalización— no es, por sí sola,
garantía de una fiscalización diferencial fuerte; el mecanismo de esta tesis
lo muestra de forma directa. Segundo, los efectos de reporte son un reto
central, y con frecuencia subestimado, para evaluar políticas de seguridad
vial con datos administrativos: el hecho de que el efecto aparente sobre
víctimas totales desaparezca al usar fallecidos —y de que ambos coincidan en
tiempo con el repunte de multas— es exactamente el patrón que produciría un
cambio en la probabilidad de registro sin un cambio real en la
siniestralidad, y cualquier evaluación de este tipo de política debería
distinguir explícitamente entre ambos canales, como se hace aquí
(Sección~\ref{sec:contexto_espacial}). Tercero, la literatura de disuasión
policial \citep{chalfinmccrary2017,deangelohansen2014} sugiere que
el margen relevante no es solo la fiscalización, sino su intensidad y su
persistencia; un análisis que solo observe si hubo o no una reforma, sin
medir directamente cuánto cambió la fiscalización efectiva, corre el riesgo de
atribuir a la política un efecto —o una ausencia de efecto— que en realidad
depende de qué tan bien se implementó.

Quedan al menos tres líneas de trabajo futuro. La más directa es cambiar la
unidad de análisis de colonia a segmento vial, lo que permitiría una
definición de exposición más precisa que un buffer de distancia y reduciría
el problema de que la infraestructura y la siniestralidad comparten
determinantes a nivel colonia (Sección~\ref{sec:datos_descriptivos}). Una
segunda línea es construir un grupo de control por \emph{matching} explícito
sobre las características previas al evento, en vez de depender enteramente
de los efectos fijos para absorber la diferencia de niveles entre colonias
tratadas y no tratadas. Una tercera es extender la ventana posterior del
análisis conforme se actualicen las fuentes de datos —en particular, C5 más
allá de febrero de 2024 e infracciones más allá de abril de 2023— lo que
permitiría evaluar con outcome de mecanismo eventos más recientes, como la
propia reforma de motociclistas de septiembre de 2023, que en esta tesis solo
pudo replicarse para el outcome de víctimas.
